\section{Finite rows, singular meshes, and physical defect floors}
\label{app:row-estimates}

We retain the unequal meshes $m=n+1$, the corrected highest
first-coordinate coefficient, and the finite polynomial values at every
collision. Throughout this appendix $n\geq32$, $L=\log n$, $H=1+L$,
and $H_n=\sum_{j=1}^n1/j$. The initial phase variables are
$\alpha,\beta\geq0$, $\alpha+\beta\leq\pi$; signed comparison objects
are obtained by separate absolute phases. Put
\begin{gather*}
 c_e=\cos((\alpha+\beta)/2),\qquad
 c_d=\cos((\alpha-\beta)/2),\qquad
 a=c_ec_d=\frac{\cos\alpha+\cos\beta}{2},\\
 p=\sin\alpha,\qquad q=\sin\beta,\qquad P=p+q,\qquad d=(q-p)_+.
\end{gather*}
The finite rows of the corrected polynomial $G_n$ are
\begin{gather*}
 \delta_r=(1-r)\bmod2,\quad
 u_{rk}=\frac{\delta_r\pi+\alpha+2k\pi}{n},\quad
 v_r=\frac{r\pi+\beta}{m},\\
 R_r=\sum_{k=0}^{n-1}|(-1)^{\delta_r}G_n(u_{rk},v_r)/2|.
\end{gather*}
Thus $S_n=(nm)^{-1}\sum_{r=0}^{2m-1}R_r$. The fixed exceptional set is
\begin{equation}\label{eq:six-rows}
 A_6=\{0,1,2m-2,2m-1,m-1,m\},\qquad
 J_6=\frac1{nm}\sum_{r\in A_6}R_r=J_{\rm end}+J_{\rm cen}.
\end{equation}

\subsection{The exact decomposition before estimates}

Write $\Phi_L(t)=\Phi_0(t)+2L|\sin\pi t|/\pi$, periodically, with
$\Phi_L(k)=1$ at integers. Its scalar series on $0<t<1$ is
\begin{equation}\label{eq:Phi-series}
 \Phi_0(t)=\frac{\sin\pi t}{\pi}
 \left\{2\gamma+\frac1t+\frac1{1-t}-2\log\pi
       +\sum_{k=1}^\infty
        \left(\frac1{k+t}+\frac1{k+1-t}-\frac2k\right)\right\}.
\end{equation}
For $0<v<2\pi$, set
\begin{align*}
 C(v)&=(p-q-2a\cot(v/2))/4,\qquad
 D(v)=(p+q+2a\cot(v/2))/4,\\
 \eta_-(v)&=\tfrac12+(\alpha-\beta+v)/(2\pi),\\
 \eta_+(v)&=\tfrac12+(\alpha+\beta-v)/(2\pi),\\
 A(v)&=2(|C(v)|+|D(v)|)/\pi,\\
 B(v)&=\frac{c_e\Phi_0(\eta_-(v))+c_d\Phi_0(\eta_+(v))}
                   {2\sin(v/2)}+\frac{I(C(v),D(v),v)}\pi.
\end{align*}
The row profile is $\mathcal H(L,v)=A(v)L+B(v)$. The complete weighted
continuation is used at every removable point. Put $h=\pi/m$,
$t=\beta/\pi$, $t_L=t$, $t_R=1-t$, and $v_L(z)=z$, $v_R(z)=2\pi-z$.
The ordinary indices on each folded side are exactly $k=2,\ldots,m-2$.
Let
\begin{gather*}
 U_0=c_e\Phi_0(\eta_-(0))+c_d\Phi_0(\eta_+(0)),\\
 B_0=4a/\pi^2,\quad D_0=U_0/\pi+B_0\log\pi,
 \quad \epsilon_n=\log(m/n),\quad Z_n=D_0-B_0\epsilon_n.
\end{gather*}
For $j\geq2$ define the convergent mesh constants
\begin{align}\label{eq:hg-definitions}
 h_j(t)&=-\log(j+t-\tfrac12)+\sum_{k=j}^\infty
 \left\{\frac1{k+t}-\log\frac{k+t+1/2}{k+t-1/2}\right\},\\
 g_j(t)&=-\tfrac12\log^2(j+t-\tfrac12)\notag\\
 &\quad+\sum_{k=j}^\infty\left\{\frac{\log(k+t)}{k+t}
 -\frac{\log^2(k+t+1/2)-\log^2(k+t-1/2)}2\right\}.\notag
\end{align}
The four-row reference is
\begin{equation}\label{eq:C2-reference}
 C_2=\ell_0(\beta)+\ell_1(\alpha,\beta)
                  +\ell_1(\alpha,-\beta)+\ell_2(\alpha,-\beta),
 \quad \ell_0(\beta)=\operatorname{sinc}\beta,
\end{equation}
where, with $w=r\pi+b$, $z_j=j\pi+\alpha$,
\[
 \ell_r(\alpha,b)=\frac{\sin b}{w}
       +8a\sum_{j=-r+1,-r+3,\ldots,r-1}\frac1{w^2-z_j^2}.
\]
The sinc form of the reference, displayed below, supplies its linked
continuous values when a rational term is removable. Set
$E_2=C_2+D_0[h_2(t)+h_2(1-t)]+B_0[g_2(t)+g_2(1-t)]$ and
$\Delta_2=E_\infty-E_2$. Define
\begin{align*}
 K_{\rm f}&=J_6+\frac1m\sum_{r\notin A_6}\mathcal H(L,v_r)+\Delta_2,
 &E_{\rm row}&=S_n-K_{\rm f},\\
 F_s(L,z)&=\mathcal H(L,v_s(z))
       -\frac{U_0+(4a/\pi)(L+\log z)}z-\frac d\pi\log(\pi-z).
\end{align*}
\begin{align*}
 Q_\tau[G]&=\frac1m\sum_{k=2}^{m-2}G(h(k+\tau))
                       -\frac1\pi\int_0^\pi G(z)\,dz,\\
 E_{\rm out}&=Q_t[F_L]+Q_{1-t}[F_R],&E_{\rm ret}&=J_6-C_2.
\end{align*}
The subtracted integrals are convergent integrals of their full
expressions. The exact singular remainders are
\begin{align}\label{eq:mesh-exact}
 R_{\rm mesh}&=-B_0\epsilon_n\mathcal H_2(t)-Z_n\mathcal H_n(t)
       -B_0\mathcal Q_n(t)-B_0L^2-2D_0L,\\
 R_{\log}&=\frac1m\bigl\{2(m-3)\log h+\log\Gamma(m-1-t)
       +\log\Gamma(m-2+t)\notag\\
 &\hspace{25mm}-\log\Gamma(2-t)-\log\Gamma(1+t)\bigr\}
           -2(\log\pi-1),\notag\\
 \mathcal H_j(t)&=h_j(t)+h_j(1-t),\qquad
 \mathcal Q_j(t)=g_j(t)+g_j(1-t).\notag
\end{align}
All gamma arguments are positive. With $E=S_n-M_\infty$ one obtains
the exact identities
\begin{align}\label{eq:complete-error-telescope}
 E&=E_{\rm row}+E_{\rm out}+E_{\rm ret}+R_{\rm mesh}
                    +\frac d\pi R_{\log},\\
 \delta_aE&=\delta_aE_{\rm row}+\delta_aE_{\rm out}
        +\delta_aE_{\rm ret}+\delta_aR_{\rm mesh}+\frac d\pi R_{\log},
\end{align}
where $\delta_aX=X(\alpha,\beta)-aX(0,0)$.
To verify these identities, subtract the two convergent series in
\eqref{eq:hg-definitions} at indices $2$ and $m-1=n$. Their cell
primitives telescope to the exact finite sums of $(k+t)^{-1}$ and
$\log(k+t)/(k+t)$. The relation
\[
 \frac1m\frac{U_0+(4a/\pi)(L+\log(h(k+t)))}{h(k+t)}
       =\frac{Z_n}{k+t}+B_0\frac{\log(k+t)}{k+t}
\]
retains the precise mesh displacement. The central logarithm sums by
$\Gamma(z+1)=z\Gamma(z)$. Finally, the defining regular finite-part
integral of $M_\infty$ cancels $\Delta_2=E_\infty-E_2$ and the two
lower mesh constants. This gives \eqref{eq:complete-error-telescope}
with the signs displayed. No limiting value replaces $S_n(0)$.

\subsection{The full phase allowance}

\begin{lemma}[The combined endpoint envelope]\label{lem:endpoint-envelope}
For either folded side and $0<z\leq\pi/4$,
\begin{equation}\label{eq:endpoint-envelope}
 \left|B(v_s(z))-\frac{U_0+(4a/\pi)\log z}{z}\right|
          \leq7+\frac P\pi[3+2\log(4\pi/z)].
\end{equation}
\end{lemma}
\begin{proof}
First, for $0<z\leq\pi/4$, $w\geq0$ and
$|C+w|,|D-w|\leq\epsilon$, the complete cancellation satisfies
\begin{equation}\label{eq:endpoint-I-perturbation}
 |I(C,D,z)-I(-w,w,z)|\leq\epsilon[11+8\log(4\pi/z)].
\end{equation}
To prove it, take disjoint pole neighborhoods of radius $z/8$ on the
cotangent circle. In either neighborhood the other cotangent has
magnitude at most $8/(7z)$, while the own cotangent, at distance
$y\leq z/8$, has magnitude at least $1/(2y)\geq4/z$.
If $w\geq4\epsilon$, the coefficient signs are preserved and the own
channel dominates in each neighborhood. Both cancellations select
the same nonsingular channel; their difference integrates to at most
$8\epsilon/7$ over both neighborhoods. If $w<4\epsilon$, the
perturbed coefficients are at most $5\epsilon$ and the reference
coefficients at most $4\epsilon$. Bounding each cancellation by twice
the other channel gives total integral at most
$72\epsilon/7<11\epsilon$. On the complement use the absolute-value
Lipschitz inequality
$|g(C,D)-g(-w,w)|\leq2\epsilon(|\cot x|+|\cot(x+z)|)$.
Each absolute cotangent has remaining integral at most
$-2\log\sin(z/8)\leq2\log(4\pi/z)$, proving
\eqref{eq:endpoint-I-perturbation}. No lower bound on $w$ was used.

The exact reference is $I(-w,w,z)=4w\log\sin z$, by integrating
$-2w\cot(x+z)$ on $(0,\pi/2-z)$ and $2w\cot x$ on
$(\pi/2,\pi-z)$. Take $w=a\cot(z/2)/2$ and $\epsilon=P/4$.
The actual left coefficients satisfy the hypotheses; reflection and
simultaneous sign reversal give the right coefficients. Therefore
\[
 |I_s(z)-2a\cot(z/2)\log\sin z|
                    \leq P[3+2\log(4\pi/z)].
\]
For $N_s=c_e\Phi_0(\eta_-)+c_d\Phi_0(\eta_+)$,
$\operatorname{Lip}\Phi_0\leq12$ and $|\Phi_0|\leq4$ give
$|N_s(z)-U_0|\leq12z/\pi$ and $|U_0|\leq8$.
The sine chord and $x-\sin x\leq x^3/6$ imply
\[
 \left|\frac{N_s(z)}{2\sin(z/2)}-\frac{U_0}z\right|
             \leq6+\pi z/6.
\]
Finally $0\leq1/x-\cot x\leq\pi x/6$ and
$0\leq-\log(\sin z/z)\leq z^2/5$ give
\[
 |2a\cot(z/2)\log\sin z-(4a/z)\log z|
               \leq az(1+|\log z|).
\]
Since $z(1+\log(1/z))\leq1$ on $(0,1]$,
$6+\pi^2/24+1/\pi<7$. Combining the last three bounds with
the exact definition of $B$ proves \eqref{eq:endpoint-envelope}.
The coarse scalar bounds used above follow directly from
\eqref{eq:Phi-series}, by pairing its positive tails; they also follow
from the stronger scalar estimates and their derivative proof in
Lemma~\ref{scalar:global-profile-bounds}.
\end{proof}

\begin{lemma}[Phase-dependent absolute error]\label{lem:phase-row-global}
For every $n\geq32$ and every closed signed phase $|\alpha|+|\beta|\leq\pi$,
\begin{equation}\label{eq:phase-global-bound}
 |S_n-M_\infty|\leq C_{E,\phi}+K_DP,
 \qquad C_{E,\phi}<\overline C_E.
\end{equation}
Here $K_D$ and $\overline C_E$ have the exact definitions in
\eqref{eq:argmax-global-error-norm}. The following formulas specify the
stronger phase allowance completely. Put $b_0=4/\pi^2$,
$r_0=\log(5/2)+1/18$, and
\begin{align*}
 D_a&=\frac{(52/\pi^2)H+7}{n}+\frac{20H}{\pi n^2},\qquad
 D_\phi=aD_a+10/n^2,\\
 \varepsilon&=3\pi/(2m),\qquad \ell=-\log\sin\varepsilon,\\
 \overline S&=1/2+2a/\varepsilon,\qquad
 \Gamma_\phi=p[1+\log(1+2\overline S/p)],\\
 M_A&=\max(p,q)/\pi+4a/\pi^2,\qquad V_A=2M_A,\\
 V_B&=23+\{3\Gamma_\phi+(2p+6P)\ell+26a+20+2d\log2\}/\pi,\\
 C_B&=6+(4+4a/\mathrm e)/\pi,\\
 G_\phi&=\frac{LV_A+V_B}{2m}+\frac{6LM_A}{m}
 +\frac4m\left\{7+\frac P\pi(5+2\log(2m))\right\}\\
 &\quad+\frac{4d\log\pi}{\pi m}+\frac{2C_B}{m},\\
 C_{6,\phi}&=\frac{21+21a+2L+\{(1+48p)H+24q\}/\pi}{n}
                    +\frac{317+5/12}{n^2},\\
 \overline D&=4(c_e+c_d)/\pi+ab_0\log\pi,\qquad
 \overline Z=\overline D+ab_0/n,\\
 M_\phi&=\frac{ab_0(2L+2r_0)}n
 +\frac{\overline D+ab_0(L+1/n)}{4n^2-1}
 +\frac{ab_0}{(2n-1)^2}\\
 &\quad+\frac{\overline Z+ab_0(\log(n+1/2)+2)}{4(n-1/2)^2},\\
 W&=\frac{3\log m+\log2+2\log\pi+2}{\pi m},\\
 C_{E,\phi}&=D_\phi+G_\phi+C_{6,\phi}+M_\phi+dW.
\end{align*}
Set $\Gamma_\phi=0$ when $p=0$. Every allowance is decreasing in real
$n\geq32$ at a fixed phase.
\end{lemma}

\begin{proof}
We establish the individual estimates in
\eqref{eq:complete-error-telescope} before maximizing any phase factor.
For the ordinary row reduction, the exact cotangent identity for $G_n$
gives
\begin{equation}\label{eq:ordinary-global-bound}
 |E_{\rm row}|\leq D_\phi+K_DP,
 \qquad K_D=\frac{(6+16/(3\pi))H+2/\pi}{n}.
\end{equation}
Lemma~\ref{lem:weighted-single-pole} gives the weighted pole error
$10/n^2+20aH_n/(\pi n^2)$. By
Lemma~\ref{lem:symmetric-reference-match}, the symmetric reference
matching error is
$\{[52a/\pi^2+16P/(3\pi)]H_n+7a+2P/\pi\}/n$.
The remaining ordinary perturbation is $6PH_n/n$; a direct variation
proof is supplied below. These three disjoint errors give
\eqref{eq:ordinary-global-bound}.

For the regular outer integral, the two actual centered-cell unions are
\[
 [h(3/2+t),\pi-h(3/2-t)],\qquad
 [h(5/2-t),\pi-h(1/2+t)].
\]
Their initial holes total $4h$ and their central holes total $2h$.
The variation and supremum bounds for the subtracted coefficient
$A(v_s(z))-4a/(\pi z)$ are $V_A$ and $M_A$. For the constant term
the sum of its two variations on the cell unions is at most $V_B$.
The latter follows from Lemma~\ref{lem:generic-cancellation}: the pole
regular part contributes $23$, the at most three monotone parameter
pieces contribute $3\Gamma_\phi/\pi$, separation contributes
$2p\ell/\pi$, the logarithmic coefficient products contribute
$6P\ell/\pi$, and the remaining regular terms contribute
$(26a+20+2d\log2)/\pi$. The bound
\[
 \int_0^S\log\frac{v+s}{|v-s|}\,ds
       \leq2v\left(1+\log\frac{S+v}{v}\right)
\]
is obtained by integrating separately on either side of $s=v$;
it is read as zero at $v=0$. It includes all moving root thresholds.
The two uniform regular-variation constants used here can be checked
as follows. The scalar series gives $|\Phi_0|\leq4$,
$\operatorname{Lip}(\Phi_0)\leq12$, $|\Phi_0''|<35$ on an open
period, and derivative jumps smaller than $6$. Each moving argument
traverses half a period, so $|N_s|\leq8$,
$\operatorname{Var}N_s\leq12$, and
$\operatorname{Var}N_s'\leq47/(2\pi)$, including the jump.
For $P_s^0=N_s/[2\sin(z/2)]-U_0/z$ one has
\[
 P_s^0(z)=\int_0^1N_s'(tz)dt+N_s(z)d_0(z),\qquad
 \operatorname{Var}P_s^0\leq\frac{47}{2\pi}+20(1/2-1/\pi).
\]
The sum of the two variations is $20+7/\pi<23$, and
$\sup|P_s^0|<6$. For the initial cancellation half, put
$w=a\cot(z/2)$. The coefficients of $\log\sin z$ in
$I_s-2w\log\sin z$ are
$\max(q-p,-2w)$ and $\max(-P,q-p-4w,-2w)$.
Their combined variation and supremum are at most $4P$ and $2P$;
the product rule gives $6P\ell$. Replacing $2w\log\sin z$ by
$4a\log z/z$ costs less than $10a$ per side: write the difference as
$(2a\cot(z/2)-4a/z)\log\sin z+4a\log(\sin z/z)/z$,
whose variations are at most $3a$ and $7a$.

On the central half put $e=\pi-z$ and
$s_L=q/2+a\tan(e/2)$, $s_R=|q/2-a\tan(e/2)|$.
Both lie in $[0,3/2]$ and have at most three monotone pieces.
Their $R$ terms in Lemma~\ref{lem:generic-cancellation} have combined
variation less than $21/2$: parameter motion costs at most
$3(1+\log4)<15/2$ and separation motion at most $3$.
Writing $k_s=2(s_s-p/2)_+$ gives $k_s(0)=d$,
$|k_s'|\leq2a$, $|k_s-d|\leq2a\tan(e/2)$.
The variation of
$(k_s-d)\log\sin e+d\log(\sin e/e)$ is less than $9/2$ per side,
using $\int_0^{\pi/2}-\log\sin e\,de=(\pi/2)\log2$ and
$\tan(e/2)\cot e\leq1/2$. Thus the central cancellation variation
is less than $20$, with supremum less than $4$.
The variation of $4a\log z/z$ on this central half costs less than
$3a$ per side. Subtracting $d\log(\pi-z)$ on the initial half costs
$d\log2$ per side. This gives exactly $26a+20+2d\log2$ in $V_B$.
Midpoint quadrature for a BV function costs half its variation divided
by $m$. On an initial hole the full constant-term envelope is
$7+(P/\pi)[3+2\log(4\pi/z)]$ by
Lemma~\ref{lem:endpoint-envelope}. Its integral to length $c$ is
$c[7+(P/\pi)(5+2\log(4\pi/c))]$, a concave function of $c$.
The paired initial lengths therefore cost at most the displayed $4/m$
term in $G_\phi$. The initial central-log subtraction costs
$4d\log\pi/(\pi m)$. Centrally the complete subtracted constant
term has supremum $C_B$, so those holes cost $2C_B/m$.
The total coefficient-hole cost is $6LM_A/m$. Consequently
\begin{equation}\label{eq:outer-global-bound}
 |E_{\rm out}|\leq G_\phi.
\end{equation}

The explicit finite-row comparison proved in the next subsection is
\begin{equation}\label{eq:four-rows-global-bound}
 |J_{\rm end}-C_2|
 \leq\frac{9+21a+\{(1+48p)H_n+24q\}/\pi}{n}
                   +\frac{317+5/12}{n^2},\qquad
 0\leq J_{\rm cen}\leq\frac{12+2L}{n}.
\end{equation}
It gives $|E_{\rm ret}|\leq C_{6,\phi}$.

For the singular mesh, use the exact paired upper cell endpoints
$n+c,n-c$, $|c|\leq1/2$. The reciprocal and logarithmic reciprocal
cell tails in \eqref{eq:hg-definitions} are at most
$1/(8(n\pm c)^2)$ and
$[\log(n\pm c)+2]/[8(n\pm c)^2]$, respectively; the integral
remainder formula for midpoint quadrature proves these estimates.
Inserting the cell tails before taking an absolute value gives
\[
 |R_{\rm mesh}|\leq
 \frac{B_0(2L+|\mathcal H_2|)}n
 +\frac{|D_0|+B_0(L+1/n)}{4n^2-1}
 +\frac{B_0}{(2n-1)^2}
 +\frac{|Z_n|+B_0[\log(n+1/2)+2]}{4(n-1/2)^2}.
\]
Indeed, the remaining paired logarithm is
$\log(1-c^2/n^2)$, whose modulus is at most $(4n^2-1)^{-1}$,
and each $|\log(1\pm c/n)|\leq(2n-1)^{-1}$.
The series give $|\mathcal H_2|\leq2r_0$. Since
$|\Phi_0|\leq4$, one has $|D_0|\leq\overline D$ and
$|Z_n|\leq\overline Z$; hence $|R_{\rm mesh}|\leq M_\phi$.
Applying BV midpoint quadrature to $\log(\pi-z)$ costs
$\log(2m)/m$ on the two cell unions. Its initial holes cost
$4\log\pi/m$, and its two central holes, whose lengths sum to $2h$,
cost at most $2(1+\log(m/\pi))/m$ by concavity of
$c(1-\log c)$. Thus $|R_{\log}|/\pi\leq W$.
Adding these estimates proves the first inequality in
\eqref{eq:phase-global-bound}.

For the uniform saving, let $D=D_a+10/n^2$ and let $G_g$ and
$\mathrm{Mesh}$ be the scalar functions in
\eqref{eq:argmax-global-error-norm}. Monotonicity of
$p[1+\log(1+(1+8am/(3\pi))/p)]$ in both $p,a$ gives
$\Gamma_\phi\leq1+\log(2+8m/(3\pi))$. Also $\ell\leq\log m$,
$C_B<8$, and
\[
 G_g-G_\phi\geq\frac{28L}{\pi^2m}(1-a),\qquad
 D-D_\phi=(1-a)D_a.
\]
Since $\overline D<49/15<4$ and $\overline Z<49/15+1/64<5$,
$M_\phi\leq\mathrm{Mesh}$ term by term. From
\eqref{eq:four-rows-global-bound},
\[
 \frac{67+19L}{n}+\frac{318}{n^2}-C_{6,\phi}
 \geq\frac{21(1-a)+(48H/\pi)(1-p)+(24/\pi)(1-q)}n.
\]
Consequently the old uniform sum $C_{E,\rm old}$ exceeds
$C_{E,\phi}$ by more than
\[
 \left\{\frac{52H}{\pi^2n}+\frac{28L}{\pi^2m}\right\}(1-a)
                  +\frac{48H}{\pi n}(1-p).
\]
Now $L>3$, $n/m\geq32/33$, and $L/H>3/4$. Using
$1-a\geq(1-\cos\alpha)/2$, the last expression divided by $H/n$
is at least
\[
 A(1-\cos\alpha)+B(1-\sin\alpha)
 \geq A+B-\sqrt{A^2+B^2}>3,
 \qquad A=\frac{398}{11\pi^2},\quad B=\frac{48}\pi.
\]
For the strict last step, $A>18/5$, $B>15$ imply
$2(A-3)(B-3)-9>27/5>0$. Thus
$C_{E,\phi}<C_{E,\rm old}-3H/n=\overline C_E$.
This is a correlation between distinct complete errors; it is not a
smaller standalone six-row estimate.

At fixed phase each rational log term decreases. For the entropy,
$m\Gamma_\phi'<p\leq\Gamma_\phi$; for the sine logarithm,
$m\ell'=\varepsilon\cot\varepsilon\leq1<\ell$.
Thus their quotients by $m$ decrease. Splitting the remaining mesh
terms into constants, $1/n$, and logarithms over their positive
quadratic denominators proves their decrease directly. This proves
the stated real-$n$ monotonicity.
\end{proof}

\subsection{The six actual rows}

For $(r,b)=(1,\beta),(1,-\beta),(2,-\beta)$ put
$w=r\pi+b$, $z=j\pi+\alpha$, $\mu=n/m$, $u=z/n$, $v=w/m$.
The allowed $j$ have parity $(1-r)\bmod2$ and $-\pi\leq u<\pi$;
there are exactly $n$ of them. These are respectively the original rows
$1,2m-1,2m-2$. Let
$\sigma_-=(-1)^{(r-j-1)/2}$, $\sigma_+=(-1)^{(j+r-1)/2}$ and
$\tau(y)=y\cot y$, $\tau(0)=1$. Direct expansion of the finite
corrected polynomial gives the collision-safe chart
\begin{align}\label{eq:finite-sinc-chart}
 \frac{F_j}{nm}
 &=\frac{c_e\sigma_-\operatorname{sinc}((z-\mu w)/2)\tau((u-v)/2)
       +c_d\sigma_+\operatorname{sinc}((z+\mu w)/2)\tau((u+v)/2)}
             {2m\sin(w/(2m))},\\
 f_j&=\frac{c_e\sigma_-\operatorname{sinc}((z-w)/2)
                 +c_d\sigma_+\operatorname{sinc}((z+w)/2)}w.
 \notag
\end{align}
Here the cosine weights use $(\alpha,b)$, which has the same product
$a$ as $(\alpha,\beta)$. The infinite absolute sum of $f_j$ is exactly
$\ell_r$ in \eqref{eq:C2-reference}. Indeed away from removable
products $f_j=4a/(w^2-z^2)$, and summing its signed reciprocal
partial fractions gives the stated finite expression. The sinc
formula extends the equality at every removable product. At the sole
zero-$w$ vertex, $(r,b,\alpha)=(1,-\pi,0)$, the linked reference has
$f_0=1$, $f_j=0$ otherwise, and $\ell_r=1$.

\begin{proof}[Proof of \eqref{eq:four-rows-global-bound}]
One has $(c_e+c_d)/w\leq1$: for $w\geq\pi$ this follows from
$c_e+c_d\leq2$, and for $w=\pi-\beta<\pi$ the two weights are
$\sin((w\pm\alpha)/2)$, whose sum is at most $w$ because
$0\leq\alpha\leq w$. If $B=[2m\sin(w/(2m))]^{-1}$, the sine
Taylor bound gives $0\leq B-1/w\leq w/(20m^2)$.

First restrict to $|j|\leq4$. There are five indices for $r=1$ and
four for $r=2$. On them $|z|\leq5\pi$,
$|(u\pm v)/2|\leq7\pi/(2n)<1$,
$|\operatorname{sinc}'|\leq1/2$, and
$0\leq1-\tau(y)\leq2y^2/5$. The last inequality follows from
$\sin y-y\cos y=\int_0^y t\sin t\,dt\leq y^3/3$ and
$\sin y\geq5y/6$. Subtracting the two formulas
\eqref{eq:finite-sinc-chart} in the order prefactor, sinc displacement,
then tau error gives at most $2/(3n^2)+1/(2n)$ per near index,
plus $\frac1{10}\sum(|z|/n+v)^2$ for the whole row. For the even
five-index set $\sum z^2\leq45\pi^2$, $\sum|z|\leq13\pi$;
for the odd four-index set these are $24\pi^2$ and $8\pi$.
The two $r=1$ rows therefore each cost
$5/(2n)+121/n^2$, and the $r=2$ row costs $2/n+75/n^2$.
Their total is $7/n+317/n^2$.

For every remaining finite index, $|z|\geq4\pi\geq2w$.
Put $D=w^2-z^2$, $D_h=\mu^2w^2-z^2$ and
\[
 K=\operatorname{sinc}((|u|+v)/2)\operatorname{sinc}((|u|-v)/2),
 \qquad H_h=(1+\cos v)/(2K).
\]
The sine chord on $[0,3\pi/4]$ and $[0,\pi/2]$ gives
$K>1/6$ and $|H_h-1|\leq u^2$. The numerator identity is
\[
 F_j=\frac{2a(1+\cos v)-p\sin u+(\sin b)\sin v}
                {2(\cos u-\cos v)}.
\]
Its normalized $a$ part is $4\mu aH_h/D_h$. The $H_h-1$ piece
costs $16a/(3n^2)$ per index. The normalization and displacement
must be kept together:
\[
 \frac\mu{D_h}-\frac1D
 =\frac1m\frac{z^2+\mu w^2}{(z^2-\mu^2w^2)(z^2-w^2)}
 \leq\frac{20}{9mz^2}.
\]
The far sums obey
\[
 \sum\frac1{|z|}\leq\frac{2H_n}\pi,\qquad
 \sum\frac1{z^2}\leq\frac3{8\pi^2}<\frac1{2\pi^2},\qquad
 w\sum\frac1{z^2}<\frac1\pi.
\]
For example the worst denominators in the square sum are
$4\pi,6\pi,8\pi,\ldots$ on each side; comparison with the integral
of $x^{-2}$ proves the displayed bound. Each row's $a$ cost is at
most $6a/n$, so the three rows cost $18a/n$. The sine part costs
$8p/(n|z|)+8qw/(nz^2)$ per index, or
$(48pH_n+24q)/(\pi n)$ over the three rows.

Outside the actual half-open mesh $|z|\geq\pi n$. Each omitted
reference tail costs at most $a/n$: use
$|f_j|\leq81a/(20z^2)$ and
$\sum_{\rm outside}z^{-2}\leq2/(\pi^2n)$.
Thus all three tails cost $3a/n$.
For the closest row, Lemma~\ref{lem:closest-row-global}, using its
exact linked $\beta=0$ formula and the weighted reciprocal-pole sum, gives
\[
 \left|\frac{R_0}{nm}-\operatorname{sinc}\beta\right|
 \leq\frac{2+H_n/\pi}{m}+\frac{\pi^2}{24m^2}
 \leq\frac{2+H_n/\pi}{n}+\frac5{12n^2}.
\]
The complete row, rather than individual finite signs, is used here.
Adding the near, far, tail and closest-row estimates proves the first
inequality of \eqref{eq:four-rows-global-bound}.

For the remaining two rows $|v-\pi|\leq\pi/m$,
$\sin(v/2)>99/100$, and $|\cot(v/2)|\leq\pi/m$.
The whole weighted polynomial row estimate, obtained by applying the
finite pole sum to \eqref{eq:finite-sinc-chart} before its linked
limits, uses $|\Phi_0|\leq4$ and
$|C|+|D|\leq P/2+a\pi/m$. Its two-row sum is at most
$(12+2L)/n$. This includes $v=\pi$ by continuity of the complete
finite polynomial row; no central integral substitutes for that row.
All the previous estimates are uniform at the zero-$w$ vertex, so
the linked sinc reference gives the same final bound there.
\end{proof}

\subsection{An ordinary cancellation perturbation}

For completeness, the sharper coefficient six in
\eqref{eq:ordinary-global-bound} follows from the following elementary
variation argument. On the cotangent circle write
\begin{gather*}
 g(C,D;x)=|C\cot x+D\cot(x+e)|-|C\cot x|-|D\cot(x+e)|,\\
 f=g(-w+c,w+d)-g(-w,w),
\end{gather*}
where $0<e\leq\pi/2$ and $|c|,|d|\leq E$. Give $g$ value zero at
the two poles. Then
\begin{equation}\label{eq:ordinary-variation}
 \operatorname{Var}_{\rm per}f\leq\frac{24E}{\sin e},\qquad
 \left|\sum_{k=0}^{n-1}f(x_0+k\pi/n)-\frac n\pi\int_0^\pi f\right|
                  \leq\frac{12E}{\sin e}.
\end{equation}
The full sharp variation bound used here and in the symmetric reference
comparison is
\begin{equation}\label{eq:sharp-cancellation-variation}
 \operatorname{Var}_{\rm per}g(C,D)\leq
                \frac{8\max(|C|,|D|)}{\sin e}.
\end{equation}
To verify it, let $M=\max(|C|,|D|)$. If $CD=0$, $g=0$.
If $CD>0$, reverse both signs to make both positive. The nonzero
intervals are $(\pi/2-e,\pi/2)$ and $(\pi-e,\pi)$.
On each the minimum of the two channel magnitudes is unimodal;
its maxima are at most $M\tan(e/2)$ and $M\cot(e/2)$, respectively.
Indeed replacing both coefficients by $M$ can only increase that
minimum, and the two unweighted magnitudes cross at the midpoint.
For $g=-2\min$, the variation including each pole jump is at most
four times the corresponding maximum. This gives
$4M(\tan(e/2)+\cot(e/2))=8M/\sin e$.
If $CD<0$, reverse signs to take $C>0>D$.
The supports are $(0,\pi/2-e)$ and $(\pi/2,\pi-e)$;
their minima respectively decrease from $|D|\cot e$ to zero and
increase from zero to $C\cot e$. Including the pole jumps gives
$4(C+|D|)\cot e\leq8M/\sin e$.
The second case is empty when $e=\pi/2$.
This proves \eqref{eq:sharp-cancellation-variation} in all cases.
The centered-cell Stieltjes argument also gives its sharp quadrature
consequence $4\max(|C|,|D|)/\sin e$. If $|w|\leq E$,
\eqref{eq:sharp-cancellation-variation} applied twice to $f$ gives
the first inequality in \eqref{eq:ordinary-variation}.

If $A=|w|>E$, the coefficient signs are preserved. After a simultaneous
sign reversal let their magnitudes be $b_1,b_2\in[A-E,A+E]$. The
function $f$ vanishes where the cotangents have opposite signs and on
each remaining support interval satisfies
\begin{gather*}
 f=-2\{\min(b_1|\cot x|,b_2|\cot(x+e)|)
                   -A\min(|\cot x|,|\cot(x+e)|)\},\\
 |f|\leq2E\cot e=:M.
\end{gather*}
The two support intervals are $(0,\pi/2-e)$ and $(\pi/2,\pi-e)$.
The branch-crossing numerator is
$[(b_1-b_2)\sin(2x+e)+(b_1+b_2)\sin e]/2$ and has at most two
zeros per period. On the nonagreeing branch in each support the
derivative has at most one zero, because the positive ratio
$\sin(x+e)/\sin x$ is strictly monotone there. There are therefore
at most four interior cuts across the two supports. A path from a
pole limit in $[-M,M]$ to the zero at the other end, with $N_i$ cuts,
has variation at most its pole-limit modulus plus $2N_iM$.
Adding its pole jump costs at most $2(N_i+1)M$. The two supports
therefore cost at most $12M\leq24E/\sin e$. At $e=\pi/2$ both
supports are empty. This proves the first part of
\eqref{eq:ordinary-variation}, including double roots and flat pieces.
Centered cells of length $\pi/n$ have a Stieltjes quadrature kernel
of modulus at most $1/2$, proving its second part whenever the samples
avoid the poles.

For an ordinary Padua row take $E=P/4$ and $e$ its distance to
$\pi\mathbb Z$. If
$k_r=\min(r,2m-r-1)$ and
$l_r=m-r-1$ on the left or $r-m$ on the right, the sine chord gives
$|\sin v_r|^{-1}\leq(m/2)(k_r^{-1}+l_r^{-1})$. Each positive index
occurs at most twice, so the sum over ordinary rows is at most
$2mH_n$. The second bound in \eqref{eq:ordinary-variation}, divided
by $nm$, is exactly $6PH_n/n$.

\subsection{The one-sided actual-origin allowance}

\begin{lemma}[The actual finite origin]\label{lem:origin-lower}
Let $D$ and $\mathrm{Mesh}$ be the explicit scalar functions in
\eqref{eq:argmax-global-error-norm}. Set
\begin{align*}
 B_{\rm out}&=\frac{\pi[L+25\log(2m)+193/3]}{8m^2}+
                     \frac1{8n^2},&
 R_+&=\left[\frac{71}{n^2}-\frac{22}{45m}\right]_+,\\
 e_0&=D+\mathrm{Mesh}+B_{\rm out}+R_+.
\end{align*}
Then $E(0)=S_n(0)-M_\infty(0,L)\geq-e_0$. The function $e_0$ is
positive and decreasing for real $n\geq32$, $R_+=0$ for $n\geq148$,
and $B_{\rm out}<1/n$ for $n\geq160$. In particular,
\begin{equation}\label{eq:actual-origin-lower}
 E(0)>-\overline e,\qquad \overline e=D+\mathrm{Mesh}+1/n,
                         \qquad n\geq160.
\end{equation}
This is a lower error estimate, not an assertion $E(0)\geq0$.
\end{lemma}
\begin{proof}
At the origin $d=0$, so \eqref{eq:complete-error-telescope} gives
$E(0)=E_{\rm row}(0)+E_{\rm ret}(0)+E_{\rm out}(0)+R_{\rm mesh}(0)$.
The previous estimates give $|E_{\rm row}(0)|\leq D$ and
$|R_{\rm mesh}(0)|\leq\mathrm{Mesh}$.

The four positive endpoint samples, with multiplicity, are
$(w,z)=(\pi,0),(\pi,0),(2\pi,\pi),(2\pi,-\pi)$.
Use $D_h=\mu^2w^2-z^2$, $D'=w^2-z^2$ and the $K,H_h$ of the
finite-row proof. At the actual origin the four endpoint rows give
$-4/m$ plus their displacement and geometric terms, whereas the
two actual central rows give exactly $J_{\rm cen}(0)=2/m$.
Consequently
\[
 E_{\rm ret}(0)=-\frac2m
   +8\sum\left(\frac\mu{D_h}-\frac1{D'}\right)
   +8\mu\sum\frac{H_h-1}{D_h}.
\]
The exact partial fraction is
\[
 \frac\mu{D_h}-\frac1{D'}
 =\frac1{2m}\left\{\frac1{(\mu w-z)(w-z)}
                    +\frac1{(\mu w+z)(w+z)}\right\}.
\]
Its bracket sum over the four samples is
\[
 S_{\rm disp}=\frac1{\pi^2}
   \left\{\frac4\mu+\frac2{2\mu-1}+\frac2{3(2\mu+1)}\right\}
 \geq\frac{56}{9\pi^2}.
\]
Each summand decreases with $\mu\in(1/2,1]$, so its value at one
is a lower bound. The sine-chord estimate in these four terms gives
$|8\mu\sum(H_h-1)/D_h|\leq71/n^2$: indeed the same near
denominators used below give $\sum D_h^{-1}\leq25/54$,
$K\geq49/50$, and $|H_h-1|\leq19/n^2$, so the coefficient is
at most $8\cdot19\cdot25/54=1900/27<71$. Since $\pi^2<10$,
\begin{equation}\label{eq:retained-origin-gain}
 E_{\rm ret}(0)>\frac{22}{45m}-\frac{71}{n^2}\geq-R_+.
\end{equation}
The central origin value used here follows directly from the linked
finite cotangent vectors below: row $m$ has one value $n$, and row
$m-1$ has positive samples with exact signed sum $n$. Both parities
therefore give $2/m$ after normalization.

For the signed outer origin, put $f(z)=\Phi_0(1/2+z/(2\pi))$ and
\begin{align*}
 b_0(z)&=2\cot(z/2)-4/z,&d_0(z)&=[2\sin(z/2)]^{-1}-1/z,\\
 P_0(z)&=2[f(z)-f(0)]/z+2f(z)d_0(z),&
 Q_0(z)&=2\cot(z/2)\log\sin z-4\log z/z,\\
 F_0(z)&=(L/\pi)b_0(z)+P_0(z)+Q_0(z)/\pi.
\end{align*}
The exact origin operator is $Q_0[F_0]+Q_1[F_0]$. Its left cell
union is $[3h/2,\pi-3h/2]$ and its right union is
$[5h/2,\pi-h/2]$. All cells lie in $I=[h/2,\pi-h/2]$.
On that interval
\begin{equation}\label{eq:origin-derivative-variation}
 \operatorname{Var}_I(F_0')
 \leq V:=L/2+(25/2)\log(2m)+37/3.
\end{equation}
Here are the component estimates, to specify the enlarged cutoff.
For $c(z)=\cot z-1/z$, its reciprocal-pole series gives
$|c''|<1$ on $[0,\pi/2]$. Thus
$\operatorname{Var}(b_0'/\pi)\leq1/2$,
$0\leq d_0\leq1/5$, $0\leq d_0'\leq1/3$, and
$|d_0''|\leq5/32$. The scalar estimates $|f'|\leq2$,
$|f''|\leq1$, $f'(0)=0$ imply
$|(2(f-f(0))/z)''|\leq8/z$ and
$\int_0^\pi|(2fd_0)''|<9$. Hence
$\operatorname{Var}_I(P_0')\leq8\log(2m)+9$.
On the initial half write
$Q_0=b_0\log\sin z+g$ with
$g=4\log(\sin z/z)/z=4\int_0^1c(tz)dt$.
The $b_0''\log\sin z$ and $g''$ terms cost respectively $1$ and $3$;
the other two differentiated product terms cost
$(9/2)\int dz/z$. On the central half, with $u=\pi-z$,
$2\tan(u/2)\log\sin u$ costs at most $4+9\int du/u$
in second derivative, and $-4\log z/z$ costs $2$.
Their first derivatives match at $\pi/2$. Thus
$\operatorname{Var}_I(Q_0')\leq(27/2)\log(2m)+10$.
Dividing by $\pi>3$ proves \eqref{eq:origin-derivative-variation}.

For a centered cell of width $h$ and a BV derivative,
\[
 \left|hF(z_c)-\int_{\rm cell}F\right|
                \leq\frac{h^2}{8}\operatorname{Var}_{\rm cell}(F').
\]
Subtract the center tangent and integrate each derivative-variation
mass to see that its kernel is at most $h^2/8$.
The two cell families therefore have total normalized lower error
at least $-\pi V/(4m^2)$.
On either initial hole, $z\leq5h/2<1$,
$d_0(z)<z/16$, $P_0(z)<7z/6$, and $g\leq0$. Since $b_0\leq0$,
\[
 F_0(z)\leq\frac76z+\frac z\pi
                    \log_+\frac{\pi}{2nz}.
\]
For $c_n=\pi/(2n)$,
$\int_0^\ell z\log_+(c_n/z)dz\leq c_n^2/4$. The two initial
holes therefore cost at most $119\pi/(24m^2)+1/(8n^2)$.
The central holes have favorable sign: there
$z\geq21\pi/22$, $\sin(z/2)^{-1}<101/100$,
$\cot(z/2)<1/13$, $b_0/\pi<-7/20$, $P_0<101/100$,
and $Q_0\leq0$. Since $L>3$, $F_0<-1/25$ on their interiors.
Subtracting their integrals improves the lower estimate. Thus
\[
 E_{\rm out}(0)\geq-\frac{\pi V}{4m^2}
                         -\frac{119\pi}{24m^2}-\frac1{8n^2}
                    =-B_{\rm out}.
\]
Together with \eqref{eq:retained-origin-gain} this proves $E(0)\geq-e_0$.

All terms of $D$ and $\mathrm{Mesh}$ decrease for real $n\geq32$
by direct differentiation. Put $K=L+25\log(2m)+193/3$.
Then
\[
 (nB_{\rm out})'=
 \frac{\pi[(1-n)K+m+25n]}{8m^3}-\frac1{8n^2}<0,
\]
because $K>64$ and $64(n-1)>26n+1$. At $n=160$, the elementary
inequalities $\log160<6$, $\log322<6$, $\pi<22/7$ give
$nB_{\rm out}<[(22/7)(6+150+193/3)+1]/1280<1$.
Also $n^2/(n+1)>147$ for $n\geq148$ and $22\cdot147>71\cdot45$,
so $R_+=0$ there. On its positive support
$R_+'=-142/n^3+22/(45m^2)<0$; afterward it is zero.
This proves all monotonicity statements and \eqref{eq:actual-origin-lower}.
\end{proof}

\subsection{The matched relative telescope}

\begin{lemma}[Complete relative allowances]\label{lem:relative-telescope}
Let $x=|\alpha|$, $y=|\beta|$, $\rho=x+y$, $u=1-a$. Use the exact
$A_4,C_4,J_N,J_W,k_{Nx},k_{Ny},k_{Wx},k_{Wy}$ in
\eqref{eq:argmax-coefficients}. For every $n\geq160$,
\begin{align}\label{eq:complete-relative-bound}
 E(\phi)-E(0)&\leq C_4+k_{Nx}x+k_{Ny}y
             +u(\overline e-A_4-\rho J_N),&&\rho\leq1/2,\\
 E(\phi)-E(0)&\leq C_4+k_{Wx}x+k_{Wy}y
             +u(\overline e-A_4-\rho J_W),&&\rho\leq1.\notag
\end{align}
All domains are closed. The individual finite allowances used to prove
this statement are valid already for $n\geq32$:
\begin{align}\label{eq:relative-components}
 \delta_aE_{\rm row}&\leq aA_4+20/n^2+\rho(8H+10)/n,\\
 |\delta_aE_{\rm ret}|&\leq
  a\rho J_N+P(7/n+4/n^2),&&\rho\leq1/2,\notag\\
 |\delta_aE_{\rm ret}|&\leq
  a\rho J_W+P(11/n+4/n^2),&&\rho\leq1,\notag\\
 |\delta_aR_{\rm mesh}|&\leq\rho/(2n),&&\rho\leq1,\notag\\
 T_{\rm out}:=\delta_aE_{\rm out}+(d/\pi)R_{\log}
 &\leq\frac\rho n\left(\frac{23}3L+c_*\right)
                      +\frac{y(17L+27)}{n^2},
 &&\begin{cases}c_*=47,&\rho\leq1/2,\\c_*=50,&\rho\leq1.
 \end{cases}\notag
\end{align}
In particular the scalar logarithmic error is included once, inside
$T_{\rm out}$; the mesh estimate is its separate companion.
\end{lemma}

\begin{proof}
We first give the ordinary reference estimate on the larger domain.
There $c_e,c_d\geq7/8$, $a\geq49/64$, $P\leq\rho\leq1$ and
$2a-P\geq17/32$. Retain $J=\lfloor n/2\rfloor$ and the two
endpoint label sets $j=2,\ldots,J$ and $j=3,\ldots,J$.
For either side put $\xi=\alpha/\pi$, $s=\pm\beta/\pi$,
$\nu=j+s$, $q_\pm=(1+s\pm\xi)/2$, and
\begin{align*}
 B_j&=\psi(j+q_+)+\psi(j+q_-)-2\log\nu,\\
 U_j&=\psi(j+q_+-\nu/(2m))-\psi(j+q_+)
       +\psi(j+q_--\nu/(2m))-\psi(j+q_-)-2\log(n/m),\\
 e_j&=\frac2{\pi^2}\left\{
  \frac{\tau_j-1}\nu B_j+\frac{\tau_j}\nu U_j\right\},
 \qquad \tau_j=\frac{\pi\nu}{2m}\cot\frac{\pi\nu}{2m}.
\end{align*}
The exact endpoint reference difference is $ae_j$. On the straight
phase path, $5j/6\leq\nu\leq7j/6$, $j+q_\pm\geq j+1/3$ and
$j+q_\pm-\nu/(2m)\geq j$. The positive digamma derivative series
gives $\psi'\leq3/(2j)$ and $|\psi''|\leq2/j^2$.
Also $|B_j|\leq2/\nu$, $|U_j|\leq5/n$, and, for
$b=\xi+|s|$, $|\dot B_j|\leq4b/j$,
$|\dot U_j|\leq3b/(mj)$. The prefactors satisfy
\[
 |(\tau_j-1)/\nu|\leq4\nu/m^2,\quad
 |[(\tau_j-1)/\nu]'|\leq12/m^2,\quad
 |\tau_j/\nu|\leq1/\nu,\quad
 |(\tau_j/\nu)'|\leq3/\nu^2.
\]
Thus
\[
 |\dot e_j|\leq\frac{2b}{\pi^2}
 \left\{\frac{29}{m^2j}+\frac{19}{m^2}
                       +\frac{22}{nj^2}+\frac4{mj^2}\right\}.
\]
The two sets have fewer than $m$ labels, harmonic sum at most $2H_n$,
and square reciprocal sum at most $3/2$. Using $H_n\leq n/4$ gives
$\sum|e_j(\phi)-e_j(0)|\leq145\rho/(27n)<8\rho/n$.
The two infinite reference tails start beyond $J$. Their derivatives
are at most $6(2/\pi^2)b/j^2$, so their total costs at most
$2a\rho/n$. Lemma~\ref{lem:relative-half-line} bounds the half-line
endpoint reference defect by $2a\rho H_n/n$: its actual offsets
$(1-s\pm\xi)/2+(j+s)/(2m)$ remain in $(0,1)$ on this domain,
and the derivative of its reciprocal-pole series is bounded by the
harmonic sum. Altogether the endpoint reference increment is at most
$a\rho(2H_n+10)/n$.
The complementary central reference labels have $z\geq\pi/2-1/m$.
Their symmetric reference is bounded by $a/\sin^2(z/2)$; its
normalized phase and origin sums are at most $2a/n$ and $2/m$.
Its relative cost is therefore at most $4a/n$.
The ordinary perturbation is $6PH_n/n$, with zero origin value;
the two pole errors cost $20/n^2+40aH/(\pi n^2)$.
Adding proves the first line of \eqref{eq:relative-components}.
At collisions pass only the complete continuous row expression to a
dense collision-free sequence. A pole sample by itself is not assigned
the polynomial value.

For the retained endpoint rows the actual positive sample sets are
$\{0\},\{0\},\{-1,1\}$ in the three charts above. On $\rho\leq1$ the
full numerator is positive because $2a-P\geq17/32$, and
\[
 w\pm z\geq(17/8)k,\qquad \mu w\pm z\geq(19/10)k,
 \qquad k=r\pm j\in\{1,3\}.
\]
The first excluded index has $|z|>w>\mu w$; thus the sample lists
are complete, with no endpoint collision. The exact finite row error is
\begin{align*}
 e_r=\frac{R_r}{nm}-\ell_r
  &=-a/m+{\sin b\over w}[\tau(v/2)-1]
   +8a\sum\left(\frac\mu{D_h}-\frac1{w^2-z^2}\right)\\
  &\quad+8\mu a\sum\frac{H_h-1}{D_h}
   +2\mu\sum\frac{-p\sin u+(\sin b)\sin v}{D_hK}.
\end{align*}
The closest row has $e_0=-a/m+\operatorname{sinc}\beta
[\tau(\beta/(2m))-1]$. Subtracting the actual origin first cancels
the four $-a/m$ terms exactly. The displacement contribution becomes
$(4a/m)[S_{\rm disp}(\phi)-S_{\rm disp}(0)]$.
A reciprocal-product derivative has magnitude at most
$(25760/104329)\rho/k^3$ on the wide domain. The eight $k$ values
are six ones and two threes, so
$|S_{\rm disp}(\phi)-S_{\rm disp}(0)|\leq3\rho/2$.
Also $K\geq49/50$, $|H_h-1|\leq19/n^2$,
$|\dot H_h|\leq5\rho/n^2$,
$|\dot D_h^{-1}|\leq(20/19)\rho/D_h$, and
$\sum D_h^{-1}\leq800/1083$. The complete difference of the $H_h$
sum therefore costs at most $148a\rho/n^2$.
The sine term costs $11P/n$. The four tau terms cost $4P/n^2$,
using $\sin\beta\geq5\beta/6$ and $\beta^2\leq6q/5$.
Thus the endpoint relative bound is
$a\rho(6/n+148/n^2)+P(11/n+4/n^2)$.
On $\rho\leq1/2$, the stronger denominators are $(21/8)k$ and
$(12/5)k$, and $\sum D_h^{-1}\leq25/54$. The same derivative
calculation gives $|S_{\rm disp}(\phi)-S_{\rm disp}(0)|\leq4\rho/5$,
$78a\rho/n^2$ for the geometric difference, and $7P/n$ for the sine
term. This gives the sharper endpoint part
$a\rho(16/(5n)+78/n^2)+P(7/n+4/n^2)$.

For each actual central row use the linked cyclic vector
\[
 L_{n,k}(\theta)=\sin\theta\cot\frac{k\pi+\theta}{n},\quad
 L_{n,0}(0)=n,\quad L_{n,k}(0)=0\ (k\ne0\bmod n).
\]
These values follow from the finite polynomial
$\sin(nt)\cot t=\cos t\,U_{n-1}(\cos t)$.
The exact two row vectors are
\begin{align*}
 F_{m,k}&=\frac{c_eL_{n,k}(\theta_-)+c_dL_{n,k}(\theta_+)}{2c},\\
 &\theta_\pm=(\alpha\pm\mu\beta)/2,\qquad c=\cos(\beta/(2m)),\\
 F_{m-1,k}&=\frac{c_eL_{n,k}(\eta_-)+c_dL_{n,k-1}(\eta_+)}{2c},\\
 &\eta_\pm=(\alpha\pm\mu\beta\pm\pi/m)/2,
 \qquad c=\cos((\pi-\beta)/(2m)).
\end{align*}
Every residue is included, and the $k-1$ is a permutation.
For $|\theta|\leq11/20$ the centered nonzero residues obey
\[
 |L'_{n,k}(\theta)|\leq\frac{2n}{5|k|}
                              +\frac{11n}{50k^2}.
\]
Their sum is at most $(4/5)nH_n+(22/25)n$; the linked zero
component costs at most $n/5$. Since $H_n\geq7/2$,
$\|L_n'\|_1\leq(9/8)nH_n$. The two origin shifts for row $m-1$
are $\pm\pi/(2m)$, not zero; their norms are less than $25n/16$.
The coefficient sum satisfies
$(c_e+c_d)/(2c)\leq8a/(7c)$, with $c>99/100$.
The vector displacement in either row thus costs at most
$(2/3)a\rho nH_n$. For row $m$ the remaining coefficient error
is at most $a\rho^2/6$, multiplying $L_n(0)$.
For row $m-1$ each coefficient error is at most
$a\rho^2/13+a\beta/m^2$. The preceding origin norm bounds their
sum by $a\rho/(4n)$ after normalization. Each complete row therefore
has relative cost at most $(3/4)a\rho H_n/n$.
Thus $|\delta_aJ_{\rm cen}|\leq(3/2)a\rho H_n/n$ and
$J_{\rm cen}(0)=2/m$, proving both retained lines of
\eqref{eq:relative-components}.

The mesh identity \eqref{eq:mesh-exact} cancels its leading terms
before estimation:
\begin{align*}
 \delta_aR_{\rm mesh}={}&-B_0\epsilon_n[\mathcal H_2(t)-\mathcal H_2(0)]
 -aZ_n(0)[\mathcal H_n(t)-\mathcal H_n(0)]\\
 &-B_0[\mathcal Q_n(t)-\mathcal Q_n(0)]
 -(D_0-aD_0(0))[\mathcal H_n(t)+2L].
\end{align*}
Differentiating the convergent series gives
$|\mathcal H_2'|\leq2$,
$|\mathcal H_n'|\leq(n-1)^{-2}$,
$|\mathcal Q_n'|\leq(L-1)/(n-1)^2$, and
$|\mathcal H_n(t)+2L|<n^{-2}$.
Also $B_0\leq1/2$, $|Z_n(0)|<5$,
$|D_0-aD_0(0)|\leq5\rho/3$. Since $t\leq\rho/3$,
\[
 |\delta_aR_{\rm mesh}|\leq\rho\left\{\frac1{3n}
            +\frac{L+9}{6(n-1)^2}+\frac5{3n^2}\right\}
 <\frac\rho n\left(\frac13+\frac{13}{168}+\frac5{96}\right)
 <\frac\rho{2n}.
\]
At zero phase it is exactly zero. The signed outer bound is proved in
the next subsection. Inserting it into
\eqref{eq:complete-error-telescope}, using $P\leq\rho$, and restoring
$a=1-u$ gives the coefficients in \eqref{eq:argmax-coefficients}:
the near constant is $16/5+10+7+47+1/2=677/10$, its inverse-square
coefficient is $78+4=82$; the wide constant is
$6+10+11+50+1/2=155/2$, its inverse-square coefficient is $148+4=152$.
The harmonic coefficient is $8+3/2=19/2$ in both cases.
Finally $E(\phi)-E(0)=\delta_aE-uE(0)$ and
Lemma~\ref{lem:origin-lower} prove \eqref{eq:complete-relative-bound}.
\end{proof}

\subsection{The signed outer term, including its scalar logarithm}

Here are details for the last component of
\eqref{eq:relative-components}. They also explain why its logarithm
must not be charged a second time. Put $e=\rho$, $\delta=\alpha-\beta$,
$f(z)=\Phi_0(1/2+z/(2\pi))$, and
\begin{align*}
 N_L&=c_ef(z+\delta)+c_df(z-e),&
 N_R&=c_ef(z-\delta)+c_df(z+e),\\
 V_s&=N_s-2af(z),&V_0&=V_s(0),\\
 P_s&=(V_s(z)-V_0)/z+V_s(z)d_0(z),&
 \mathcal C_s&=I_s(\phi,z)-aI_s(0,z)-d\log(\pi-z).
\end{align*}
The exact decomposition is
\begin{align}\label{eq:signed-outer-exact}
 T_{\rm out}={}&L\{Q_t[g_L]+Q_{1-t}[g_R]\}
 +Q_t[P_L]+Q_{1-t}[P_R]
 +\{Q_t[\mathcal C_L]+Q_{1-t}[\mathcal C_R]\}/\pi\\
 &+\mathrm{Shift}+(d/\pi)R_{\log},\notag\\
 \pi g_L&=\max(q,p-Z),\qquad
 \pi g_R=\max(-q,p-Z,q-2Z),\qquad Z=2a\cot(z/2),\notag\\
 \mathrm{Shift}&=\frac am\sum_{k=2}^{m-2}
 \{F_0(h(k+t))+F_0(h(k+1-t))-F_0(hk)-F_0(h(k+1))\}.\notag
\end{align}
First, $v_L=g_L-q/\pi=(p-q-Z)_+/\pi$ and
$v_R=g_R+q/\pi=\max(0,p+q-Z,2q-2Z)/\pi$ are nonnegative and
nondecreasing. For any such $G$, the disjoint right rectangles
$[h(k+\tau),h(k+\tau+1)]$, $2\leq k\leq m-2$, lie in $[0,\pi]$
and show $Q_\tau[G]\leq0$. Because $Q_\tau[1]=-3/m$,
the opposite constants cancel and
$Q_t[g_L]+Q_{1-t}[g_R]\leq0$. Thus the entire $L$ coefficient
in \eqref{eq:signed-outer-exact} has a favorable sign at every degree.

The scalar estimates $0<f\leq1$, $\operatorname{Lip}(f)<2$,
$|f''|<1$ on an open smooth period, $f'(0)=0$, and
$\operatorname{Var}_{\rm per}f'<4$ are used here. The last bound
includes the negative cusp atom: the full measure $df'$ has total
mass zero and positive density at most $3/\pi^2$, hence total
variation at most $12/\pi<4$. The endpoint-cancelled scalar integral
\[
 f(z)=f(0)\cos(z/2)+2\int_0^{1/2}
                \frac{\cos(zt)-\cos(z/2)}{\cos\pi t}\,dt
\]
also gives $|f'''|<13/20$ on $(-\pi,\pi)$: the differentiated
numerator vanishes at $t=1/2$, its $t$ derivative is at most
$3/4+\pi/8<23/20$, and $\cos\pi t\geq2(1/2-t)$.

Set $r=c_e+c_d-2a$, so $0\leq r\leq\rho^2/4$.
The following table lists bounds for each side. In its first two
columns the interval is $[0,\pi/2]$; in its last two it is
$[\pi/2,\pi]$. Each entry is multiplied by $\rho$.
\begin{center}
\begin{tabular}{c|cccc}
 domain&$\operatorname{Var}P_s$&$\sup|P_s|$
       &$\operatorname{Var}P_s$&$\sup|P_s|$\\ \hline
 $\rho\leq1/2$&$537/200$&$123/40$&$731/96$&$337/96$\\
 $\rho\leq1$&$289/100$&$67/20$&$395/48$&$193/48$
\end{tabular}
\end{center}
To verify the wide row, translations and $r f$ give
$|V_s|\leq17\rho/4$, $|V_0|\leq5\rho/4$,
$|V_s'|\leq5\rho/2$, and $|V_s''|\leq31\rho/20$ on the initial
half. Preserve the quotient before differentiation:
\[
 \left((V_s-V_0)/z\right)'
             =z^{-2}\int_0^z wV_s''(w)\,dw.
\]
Its variation is at most $31\rho/25$, while that of $V_sd_0$ is
at most $33\rho/20$. This proves the first wide entry and its
supremum bound. Centrally the measure translation estimate
$\operatorname{Var}[f(\cdot+b)-f]\leq4|b|$ yields
$\operatorname{Var}V_s\leq9\rho$. Use the exact identity
$P_s=V_s/[2\sin(z/2)]-V_0/z$. The two weights have
supremum/variation at most $(3/4,1/4)$ and $(2/3,1/3)$,
giving the last two entries. The half-radius row follows by retaining
$r\leq\rho/8$, which gives
$|V_s|\leq33\rho/8$, $|V_0|\leq5\rho/8$,
$|V_s'|\leq9\rho/4$, $|V_s''|\leq57\rho/40$ and
$\operatorname{Var}V_s\leq17\rho/2$ centrally.
Midpoint BV error, the initial holes of total $4h$ and the central
holes of total $2h$ give respectively
\[
 |Q_t[P_L]+Q_{1-t}[P_R]|
 \leq\begin{cases}30\rho/m,&\rho\leq1/2,\\33\rho/m,&\rho\leq1.
 \end{cases}
\]
For example the wide coefficient is
$289/100+395/48+4(67/20)+2(193/48)=39073/1200<33$.

We retain the complete cancellation charge. On the initial half
$s_R=a\cot(z/2)-q/2$ satisfies
$s_R\sin z\geq a-q/2\geq p/2+17/64$; both channels are root-free.
Their exact increments are $q\log\sin z$ and $-q\log\sin z$.
On the central half put $e=\pi-z$, $s_0=a\tan(e/2)$ and
$s_L=s_0+q/2$, $s_R=|s_0-q/2|$. The exact expression is
\[
 \mathcal C_s=R(p/2,s_s,e)+b_s(e)\log\sin e+d\log(\sin e/e),
\]
where $b_L=\max(q-p,-2s_0)-d$ and
$b_R=\max(-p-q,q-p-4s_0,-2s_0)-d$.
The $b$ functions decrease from zero, and their combined supremum and
variation are at most $3\rho$. The two $s$ curves have at most three
monotone pieces with $s\leq3/2$.
Let $\varepsilon=h/2$, $\ell=-\log\sin\varepsilon$.
If $\rho\leq\varepsilon^2/8$, then
$[a\tan(e/2)-q/2]\sin e\geq7e^2/80\geq\rho/2\geq p/2$
for $e\geq\varepsilon$; there are no central roots on the cells.
Otherwise the parameter integral is bounded by
$p[1+\log(4/p)]\leq\rho[11/2+2\ell]$,
using $p\log(\rho/p)\leq\rho/\mathrm e$.
Lemma~\ref{lem:generic-cancellation} in these two alternatives yields
full two-side variation at most $\rho(39/2+16\ell)$.
The initial-hole envelope is
$\rho[\log(\pi/(2z))+\log\pi]$; the central-hole envelope is
$\rho[7+4(-\log\sin e)]$. Both are integrable. The midpoint,
initial-hole, and central-hole charges are respectively at most
\[
 \frac\rho{\pi m}(39/4+8\log m),\quad
 \frac{4\rho}{\pi m}(3+\log m),\quad
 \frac\rho{\pi m}(22+8\log m).
\]
Thus the complete cancellation cost is
$\rho[175/4+20\log m]/(\pi m)$, including root and coefficient
switches. The scalar logarithm costs $dW$ as already proved.

The remaining shift is bounded by
$|\mathrm{Shift}|\leq\beta(17L+27)/m^2$.
For specificity, the first differences of its finite sum telescope
to the two boundary differences, and the interior terms become
disjoint second differences. On $[2h,\pi-h]$ the fixed profile obeys
\[
 \sup|F_0'|\leq(4/3)L+(2/3)\log m+20/3,\qquad
 \operatorname{Var}(F_0')\leq L/2+(25/2)\log m+37/3.
\]
The two boundary differences and the interior derivative-variation
measure therefore cost
$\beta[2\sup|F_0'|+\operatorname{Var}(F_0')]/m^2$,
which is less than $\beta(17L+27)/m^2$ after
$\log m\leq L+1/32$. These estimates follow by the same
$b_0,P_0,Q_0$ differentiation used in
\eqref{eq:origin-derivative-variation}, now with cutoff $2h$.
All arguments remain in that interval for $0\leq t<1/3$.

Adding every term in \eqref{eq:signed-outer-exact} gives
\[
 T_{\rm out}\leq\frac\rho m
 \left\{c_P+\frac{23\log m+175/4+\log2+2\log\pi+2}\pi\right\}
       +\frac{\beta(17L+27)}{m^2},
\]
where $c_P=30$ or $33$ in the two domains. Using $\pi>3$,
$\log2<7/10$, $\log\pi<23/20$, $\log m\leq L+1/32$ proves
the constants $47$ and $50$ in \eqref{eq:relative-components}.
Every exact difference vanishes at the origin. All scalar right sides
decrease for real $n\geq32$ by differentiation.

\subsection{The physical projected floors}

\begin{lemma}[Three center classes]\label{lem:projected-floors}
Let $n\geq64$, and let an applicable horizontal or vertical projected
defect belong to an interior center index. For actual phase radius
$\rho\leq r$, put $a_r=1-r^2/4$, $p_\pi=22/7$, and
\begin{align*}
 \eta&=a_r[1-(4p_\pi+r)/n]-r/2,\\
 F_n(r)&=\frac{2\eta}{p_\pi+r}\frac n{n+1}(1/4-1/n),\\
 F_s(r)&=\frac{35}{286}(1-6/n)^2(a_r-r/2),\\
 A_0(r)&=97a_r/100-r/2,\\
 I_h(r)&=2303a_r/5000-471r/1250,\\
 R_h(r)&=4837a_r/1875-5r/3,\\
 I_w(r)&=5369a_r/15000-7379r/30000,\\
 R_w(r)&=37621a_r/15000-3r/2,\\
 F_h(r)&=\frac{539}{1936}I_h(r)-\frac{140A_0(r)}{157n}
               -\frac{400R_h(r)}{157n}-\frac{800A_0(r)}{3n^2},\\
 F_w(r)&=\frac{539}{1936}I_w(r)-\frac{140A_0(r)}{157n}
               -\frac{400R_w(r)}{157n}-\frac{800A_0(r)}{3n^2}.
\end{align*}
Each applicable projection is at least
\begin{equation}\label{eq:three-class-floors}
 \min(F_n,F_s,F_h)\quad(0\leq r\leq21/40),\qquad
 \min(F_n,F_s,F_w)\quad(0\leq r\leq\pi/4).
\end{equation}
The minimum concerns three exhaustive classes of center positions;
the separate terms are not individual bounds for every center.
The full physical defect is at least either eligible projection, so
these formulas apply to every noncorner odd center with an interior
index. They are never applied to an actual corner. At a fixed legal
$r$, each floor increases with real $n\geq64$.
\end{lemma}

\begin{proof}
Use the exact physical defect projection in the geometric appendix.
After folding its reflected preimages, an interior physical row has
normalization $1/(nm)$; both opposite periodic branches are retained.
The original endpoint half weights stay in the full definition, and
omitted terms are nonnegative. For regular denominators the common
numerator in a selected pair is
\[
 A(u,v)=a(1+\cos v)-(\sin\alpha/2)\sin u
                         +(\sin\beta/2)\sin v.
\]
Here $a\geq a_r$ and
$(|\sin\alpha|+|\sin\beta|)/2\leq r/2$. In the horizontal direction
the baseline remains $1+\cos v$ at its fixed actual second coordinate;
the two coordinate meshes are not interchanged.
Let $M,L$ denote the adjacent physical and periodic meshes:
$(m,n)$ vertically, $(n,m)$ horizontally. They satisfy $M,L\geq n$
and $M/L\leq65/64$. If the center angle is at distance $\gamma$
from its nearest endpoint, the two cosine targets in physical row
$b=l\pi/M$ are $z_0\pm D$ with
$z_0=\cos\gamma\cos b$, $D=\sin\gamma\sin b$.
A selected periodic cosine strictly between the targets has opposite
denominator signs; a positive numerator lower $A$ gives pair defect
at least $2A/(D+|z-z_0|)$.

If the center-index distance from an endpoint is $1\leq s\leq4$,
choose physical rows $M/4\leq l\leq M/2$ and periodic offsets
$-(s-1),-(s-3),\ldots,s-1$ in both branches. The displacement estimate
\[
 |M(w-b)|\leq(s-1)\pi+71\pi/128
                         +(65/64)|\text{periodic phase}|
\]
puts every selected angle strictly between the targets: the sufficient
condition is
$|\text{physical phase}|+(65/64)|\text{periodic phase}|<57\pi/128$,
and its left side is at most $65\pi/256$ when $r\leq\pi/4$.
The parity indices remain inside the physical range
($l+q\geq13$, $l+q\leq M/2+3<L$); the two branches are distinct.
The target half-width is $t\leq(s\pi+r)/M<\pi/4$.
Both actual cosine targets, and the horizontal fixed baseline,
are at least $-t$. Their numerator is therefore at least $\eta>1/4$:
indeed $r<11/14$, $a_r>663/784$ and
$(4p_\pi+r)/n<187/896$ imply
$\eta>194099/702464$.
Each pair costs at least $\eta/t$. There are $s$ offsets, two
branches and at least $M/4-1$ physical rows. Hence
\[
 Q\geq\frac{2s\eta}{s\pi+r}\frac ML(1/4-1/M)\geq F_n(r).
\]

For remaining center indices $s\geq5$, one has $\gamma>4\pi/L$.
If $\gamma<\pi/6$, use rows $b\in[\pi/6,\pi/3]$ and the two open
periodic arcs $(b-\gamma,b+\gamma)$ and its reflection. Their combined
count is at least $2L\gamma/\pi-2$ and their numerator is at least
$A_\gamma=a_r[1+\cos(\pi/3+\gamma)]-r/2>0$.
The physical-row count is at least $M/6-1$, and each pair contributes
$2A_\gamma/(\sqrt3\gamma)$. Thus
\[
 Q\geq\frac4{\pi\sqrt3}(1/6-1/M)h(\gamma),\quad
 h(g)=A_g(1-\pi/(Lg)).
\]
The function $A$ is positive, decreasing and concave on
$[4\pi/L,\pi/6]$. Direct differentiation gives $h''<0$,
so its minimum is at an endpoint. The lower endpoint minus the upper
is bounded below by
$[1/8+(6-3\pi)/L]a_r+[1/8-3/L]r>0$.
Consequently $h\geq(a_r-r/2)(1-6/L)$.
Using $\pi<22/7$ and $\sqrt3<26/15$ gives $Q\geq F_s(r)$.
This keeps the count and numerator at the same $\gamma$.

For $\gamma\geq\pi/6$, take a positive physical band
$[a_0,b_0]$, where $a_0=\pi/100$ and either $b_0=\pi/6$ or
$b_0=\pi/4$. Write $C=a_r-r/2$,
$A(b)=C-a_r\sin b$, and $T_0=4\log2$. The two open-arc
layer-cake count, followed by the finite physical mesh sum, gives
\begin{equation}\label{eq:positive-band-floor}
 Q\geq\frac{T_0}{\pi^2}I-\frac{T_0A_{0,\rm ex}}{\pi n}
       -\frac{8R_{\rm ex}}{\pi n}
       -\frac{8A_{0,\rm ex}}{n^2\sin a_0},
\end{equation}
where
\begin{gather*}
 I=C(b_0-a_0)+a_r(\cos b_0-\cos a_0),\qquad
 A_{0,\rm ex}=C-a_r\sin a_0,\\
 R_{\rm ex}=C\log\frac{\tan(b_0/2)}{\tan(a_0/2)}-a_r(b_0-a_0).
\end{gather*}
To see the losses in \eqref{eq:positive-band-floor}, integrate the
two periodic-arc superlevel counts; the centered integral yields
$T_0\int A/\pi^2$. The two open-arc count losses are bounded by
$8\int A/\sin b\,db/(\pi n)$. Replacing the finite physical sum by
its integral loses its first value $T_0A_{0,\rm ex}/(\pi n)$
and the product of the two mesh endpoint losses
$8A_{0,\rm ex}/(n^2\sin a_0)$. The monotonicity of $A$ and
$A/\sin b$ makes each replacement a lower estimate.

On the wider radius domain choose $b_0=\pi/6$.
Then $A(b)\geq(a_r-r)/2>47/1568>0$.
The elementary bounds
$\sin(\pi/100)>3/100$, $\cos(\pi/6)>433/500$,
$b_0-a_0>7379/15000$, and
$\log[\tan(\pi/12)/\tan(\pi/200)]<3$
give $I>I_w$, $A_{0,\rm ex}<A_0$, $R_{\rm ex}<R_w$.
The positive series for $\log2$ gives $11/4<T_0<14/5$;
$157/50<\pi<22/7$ inserted with the correct signs in
\eqref{eq:positive-band-floor} gives $F_w$ exactly.

For $b_0=\pi/4$, the same positivity holds through $r=21/40$:
\[
 A(b)>\frac{29}{100}(1-r^2/4)-r/2
       \geq\frac{4811}{640000}>0.
\]
The endpoint comparison uses $\sin(\pi/4)<71/100$.
The inequalities $b_0-a_0=6\pi/25$,
$\cos(\pi/4)>707/1000$ and the logarithmic ratio bound $<10/3$
now give $I>I_h$, $R_{\rm ex}<R_h$, hence $F_h$ by the same full
finite formula. This proves its stated $21/40$ guard directly.
The nearest, small-angle and middle-angle classes exhaust both
eligible projections, including all seams and both parities.
At a fixed legal $r$, $\eta$ and the positive nearest factors increase
with $n$, as does $F_s$; the positive loss coefficients in $F_h,F_w$
are multiplied by decreasing $1/n$ and $1/n^2$. Their minimum
therefore increases. This proves all assertions.
\end{proof}
